
\subsection*{2.3.2}
\phantomsection
\label{sga3:VIIB:2.3.2}
\addcontentsline{toc}{subsection}{\SGARefLink{sga3:VIIB:2.3.2:prop}{2.3.2}}

Let us now suppose, for simplicity, that the ring \(k\) is Artinian.
When \(G\) is topologically flat over \(k\), the algebra
\[
  \mathbf H(G)=\mathbf H(G)(k)
\]
may be characterized by a universal property due to Cartier. Recall
\textup{(cf.\ \SGARefLink{sga3:VIIA:1.2.1}{1.2.1})} that, if \(U\) is a \(k\)-module, we write
\(\mathbf W(U)\) for the functor that assigns to every finite-length
\(k\)-algebra \(C\) the \(C\)-module
\[
  U\otimes_k C.
\]
\footnote{\label{sga3:VIIB:nde:88}The editors added the preceding
sentence and, in what follows, wrote \(\mathbf W(U)^\times\) in
place of \(U_m\). Recall also \textup{(cf.\ \SGARefLink{sga3:VIIA:1.2.1}{1.2.1})} that a
``\(k\)-functor'' is a covariant functor from
\(\operatorname{Alf}/k\) to \((\mathbf{Sets})\).}
If \(U\) is a \(k\)-algebra \textup{(associative and with a unit
element)}, then so is \(U\otimes_k C\). We shall write
\(\mathbf W(U)^\times\) for the \(k\)-group functor that assigns to
every \(C\in\operatorname{Ob}(\operatorname{Alf}/k)\) the
multiplicative group of invertible elements of the algebra
\(U\otimes_k C\):
\[
  \mathbf W(U)^\times(C)=(U\otimes_k C)^\times.
\]

Moreover, identify \(G\) with the \(k\)-group functor
\[
  C\longmapsto
  \operatorname{Hom}_{\operatorname{Vaf}/k}
    \bigl(\operatorname{Spf}(C),G\bigr),
\]
and write
\[
  \operatorname{Hom}_{k\text{-}\operatorname{Gr.}}
    \bigl(G,\mathbf W(U)^\times\bigr)
\]
for the set of homomorphisms of \(k\)-group functors from \(G\) to
\(\mathbf W(U)^\times\). We have the following.

\medskip
\noindent\textbf{Proposition.}
\phantomsection
\label{sga3:VIIB:2.3.2:prop}
{\itshape Let \(k\) be an Artinian ring. For every formal group \(G\)
topologically flat over \(k\), and every \(k\)-algebra \(U\), there
is a canonical isomorphism
\[
  \operatorname{Hom}_{k\text{-}\operatorname{Gr.}}
    \bigl(G,\mathbf W(U)^\times\bigr)
  \xrightarrow{\ \sim\ }
  \operatorname{Hom}_{k\text{-}\operatorname{Alg.}}
    \bigl(\mathbf H(G),U\bigr).
\]}

Let \(A\) be the affine algebra of \(G\). By hypothesis, \(A\) is a
projective object of \(\operatorname{PC}(k)\), and
\[
  \mathbf H(G)=A^\dagger.
\]
For every object \(P\) of \(\operatorname{PC}(k)\), write
\[
  U\widehat\otimes_k P
\]
for the inverse limit of the \(k\)-modules
\[
  U\otimes_k(P/N),
\]
where \(N\) runs through the open submodules of \(P\). There are
linear maps
\[
  U\otimes_k(P/N)
  \longrightarrow
  \operatorname{Hom}_k\bigl((P/N)^*,U\bigr)
\]
that send \(u\otimes x\) to the \(k\)-linear map
\[
  f\longmapsto f(x)u
\]
and form a filtered inverse system. Passing to the inverse limit
therefore gives a morphism
\begin{equation}
  U\widehat\otimes_kP
  \xrightarrow{\ \psi_P\ }
  \operatorname{Hom}_k(P^\dagger,U).
  \tag{1}
  \label{sga3:VIIB:2.3.2:eq1}
\end{equation}

When \(P=k\), the morphism \(\psi_k\) is plainly an isomorphism.
Moreover, the two sides of \textup{\SGARefLink{sga3:VIIB:2.3.2:eq1}{(1)}}, considered as functors of
\(P\), commute with infinite products \textup{(every product being
a filtered inverse limit of finite products)}. Thus \textup{\SGARefLink{sga3:VIIB:2.3.2:eq1}{(1)}}
is an isomorphism when \(P\) is a product of copies of \(k\), and
then when \(P\) is a projective object of \(\operatorname{PC}(k)\)
\textup{(the two sides of \textup{\SGARefLink{sga3:VIIB:2.3.2:eq1}{(1)}} commute with finite direct
sums)}.

Now let
\[
  \operatorname{Hom}_{\mathcal F}\bigl(G,\mathbf W(U)\bigr)
\]
denote the set of morphisms of \(k\)-functors from \(G\) to
\(\mathbf W(U)\). Since
\[
  G=\operatorname{Spf}(A)
   =\varinjlim_{\mathfrak l}\operatorname{Spf}(A/\mathfrak l),
\]
where \(\mathfrak l\) runs through the open ideals of \(A\), there
are canonical isomorphisms
\[
\begin{aligned}
  \operatorname{Hom}_{\mathcal F}\bigl(G,\mathbf W(U)\bigr)
  &=
  \varprojlim_{\mathfrak l}
  \operatorname{Hom}_{\mathcal F}
    \bigl(\operatorname{Spf}(A/\mathfrak l),\mathbf W(U)\bigr)
  \\
  &=
  \varprojlim_{\mathfrak l}
  U\otimes_k(A/\mathfrak l)
  =
  U\widehat\otimes_kA.
\end{aligned}
\]
The preceding argument therefore gives a canonical isomorphism
\begin{equation}
  \operatorname{Hom}_{\mathcal F}\bigl(G,\mathbf W(U)\bigr)
  =
  U\widehat\otimes_kA
  \xrightarrow{\ \psi_A\ }
  \operatorname{Hom}_k\bigl(\mathbf H(G),U\bigr).
  \tag{2}
  \label{sga3:VIIB:2.3.2:eq2}
\end{equation}

For every finite-length \(k\)-algebra \(C\), multiplication makes
\(U\otimes_kC\) a monoid with a unit element, and every homomorphism
of monoids with unit
\[
  G(C)\longrightarrow U\otimes_kC
\]
is necessarily a group homomorphism
\[
  G(C)\longrightarrow(U\otimes_kC)^\times.
\]
It follows that
\[
  \operatorname{Hom}_{k\text{-}\operatorname{Gr.}}
    \bigl(G,\mathbf W(U)^\times\bigr)
\]
is the part of
\(\operatorname{Hom}_{\mathcal F}(G,\mathbf W(U))\) consisting of
morphisms of \(k\)-functors into monoids with a unit element.

It remains to show that these morphisms correspond to the
\(k\)-linear maps from \(\mathbf H(G)\) to \(U\) that preserve
multiplication and unit elements.
\footnote{\label{sga3:VIIB:nde:89}\textit{Editors' note.} The original stated that this
``follows from the functoriality of \(\psi_A\).'' The editors gave
the details that follow.}
To simplify the notation, write
\[
  H=\mathbf H(G)=A^\dagger
\]
and write \(\widehat\otimes\) instead of
\(\widehat\otimes_k\). Denote by
\[
  \Delta_A,\ m_A,\ \varepsilon_A
  \qquad\text{and}\qquad
  \Delta_H,\ m_H,\ \varepsilon_H
\]
the comultiplication, multiplication, and augmentation of \(A\)
and \(H\), respectively. Let
\[
  \theta\in
  \operatorname{Hom}_{\mathcal F}\bigl(G,\mathbf W(U)\bigr),
\]
let \(\gamma\) be its image in \(U\widehat\otimes A\), and let
\[
  \phi:H\longrightarrow U
\]
be the associated \(k\)-linear map. The morphism \(\theta\) sends
the unit section \(s\in G(k)\) to an element \(u\) of \(U\). Since
\(s\) corresponds to the augmentation
\[
  \varepsilon:A\longrightarrow k,
\]
which is the unit element \(1_H\) of \(H\), we have
\[
  \theta(s)=1_U
  \quad\Longleftrightarrow\quad
  \phi(1_H)=1_U.
\]

The morphism
\[
  \theta\circ m_G:G\times G\longrightarrow\mathbf W(U)
\]
corresponds to the element
\[
  (\operatorname{id}_U\widehat\otimes\Delta_A)(\gamma),
\]
which, by duality, corresponds to the map
\[
  \phi\circ m_H:H\otimes H\longrightarrow U.
\]

On the other hand, the morphism
\[
  \theta\circ\operatorname{pr}_1:
  G\times G\longrightarrow\mathbf W(U)
\]
corresponds to the element
\[
  \gamma\widehat\otimes1_A
  \in
  U\widehat\otimes A\widehat\otimes A,
\]
which, by duality, corresponds to
\[
  \phi\circ(\operatorname{id}_H\otimes\varepsilon_H):
  H\otimes H\longrightarrow U.
\]
Similarly, \(\theta\circ\operatorname{pr}_2\) corresponds to the
element
\[
  \tau(\gamma\widehat\otimes1_A)
  \in
  U\widehat\otimes A\widehat\otimes A,
\]
where
\[
  \tau(u\widehat\otimes a\widehat\otimes b)
  =
  u\widehat\otimes b\widehat\otimes a,
\]
and this element corresponds by duality to
\[
  \phi\circ(\varepsilon_H\otimes\operatorname{id}_H):
  H\otimes H\longrightarrow U.
\]

Finally, the multiplication map
\[
  \mu=m_{U\widehat\otimes A\widehat\otimes A}
\]
given by
\[
\begin{aligned}
  &(U\widehat\otimes A\widehat\otimes A)
   \times
   (U\widehat\otimes A\widehat\otimes A)
  \longrightarrow
  U\widehat\otimes A\widehat\otimes A,
  \\
  &\bigl(
    u\widehat\otimes a_1\widehat\otimes a_2,\,
    u'\widehat\otimes a_3\widehat\otimes a_4
  \bigr)
  \longmapsto
  uu'\widehat\otimes a_1a_3\widehat\otimes a_2a_4
\end{aligned}
\]
may be viewed as the composite of the endomorphism
\(\sigma_{23}\) of
\[
  (U\otimes U)\widehat\otimes A^{\widehat\otimes4},
\]
which ``exchanges the factors \(a_2\) and \(a_3\),'' and the map
\[
  (U\otimes U)
  \widehat\otimes A^{\widehat\otimes2}
  \widehat\otimes A^{\widehat\otimes2}
  \xrightarrow{\ m_U\widehat\otimes m_A\widehat\otimes m_A\ }
  U\widehat\otimes A\widehat\otimes A.
\]

It follows that the map
\[
  \mu\circ
  \bigl(\theta\circ\operatorname{pr}_1,\,
        \theta\circ\operatorname{pr}_2\bigr):
  G\times G\longrightarrow\mathbf W(U)
\]
corresponds to the composite map \(\beta\) below:
\begin{center}
  \phantomsection
  \label{sga3:VIIB:2.3.2:beta-diagram}
  \begin{tikzcd}[row sep=2.7em,column sep=3em]
{} & H\otimes H
     \arrow[r,"\beta"]
     \arrow[dl,"{\operatorname{id}_H\otimes\sigma_{23}\otimes
                   \operatorname{id}_H}"']
   & U \\
H\otimes H\otimes H\otimes H
  \arrow[dr,"{(\operatorname{id}_H\otimes\varepsilon_H)
              \otimes(\varepsilon_H\otimes\operatorname{id}_H)}"']
  & {} & U\otimes U \arrow[u,"m_U"'] \\
{} & H\otimes H \arrow[ur,"\phi\otimes\phi"'] & {.}
\end{tikzcd}

\end{center}

Finally, \(\theta\) is compatible with the laws of \(G\) and
\(\mathbf W(U)^\times\) if and only if
\[
  \theta\circ m_G
  =
  \mu\circ
  \bigl(\theta\circ\operatorname{pr}_1,\,
        \theta\circ\operatorname{pr}_2\bigr),
\]
\noindent\textit{Translator's note.}\enspace
The French re-edition prints \(\phi\circ\operatorname{pr}_i\) in
these two ordered pairs. Since
\(\operatorname{pr}_i:G\times G\to G\) whereas
\(\phi:H\to U\), the required expression is the
type-correct \(\theta\circ\operatorname{pr}_i\).

which, by the preceding argument, is equivalent to
\[
  \phi\circ m_H=\beta.
\]
But plainly
\[
  (\phi\circ m_H)(x\otimes y)=\phi(xy),
\]
whereas
\[
  \beta(x\otimes y)=\phi(x)\phi(y).
\]
